\documentclass[a4paper,10pt]{article}

\usepackage[margin=2cm]{geometry}
\usepackage{graphicx}
\usepackage{hyperref}
\usepackage[all]{hypcap}
\usepackage{tabu}
\usepackage[title,titletoc,toc]{appendix}
\usepackage[english]{babel}
\usepackage{fontspec}
\usepackage[style=authoryear,backend=biber,sorting=nyt,dashed=false,urldate=long,abbreviate=false]{biblatex}
\usepackage{float}
\usepackage{fancyhdr}
\usepackage{microtype}

\addbibresource{references.bib}

\setlength{\headheight}{15.2pt}
\pagestyle{fancy}
\lhead{}
\chead{}
\rhead{\bfseries Software Architecture Specification}
\lfoot{Team Echo}
\cfoot{COS301 Software Engineering}
\rfoot{Page \thepage}
\renewcommand{\headrulewidth}{0.4pt}
\renewcommand{\footrulewidth}{0.4pt}

\setlength{\parindent}{0pt}
\setlength{\parskip}{1ex plus 0.5ex minus 0.2ex}

\frenchspacing

\title{\includegraphics[width=12cm]{Eeufeeslogo.jpg} \\
       Software Architecture \\
       Research Paper Management System \\
       \vspace{0.5cm}
       University of Pretoria \\
       \vspace{1.0cm}
       }

\date{} 
\author{Team Echo\\
	\vspace{0.5cm} \\
	\begin{tabu} to \textwidth { X[l] X[l]}
		\hline
		\textbf{Surname, First Name (Initial)}	& \textbf{Student Number}	\\ \hline \hline
		Bode, Elizabeth (EF)			& 14310156		\\ \hline
		Bondjobo, Jocelyn (JM)		& 13232852		\\ \hline
		Broekman, Andrew (A)		& 11089777		\\ \hline
		Loreggian, Fabio (FR)			& 14040426		\\ \hline
		Schutte, Gerome (GC)		& 12031519		\\ \hline
		Sefako, Motsitsiripe (MG)		& 12231097		\\ \hline
		Singh, Emilio (E)			& 14006512		\\ \hline
		\hline
	\end{tabu}}

\DefineBibliographyStrings{english}{%
urlseen = {\mbox{Accessed} on},
urlfrom = {[Online]},
in = {In},
}
\DeclareFieldFormat{url}{\bibstring{urlfrom}\space\url{#1}}
\DeclareFieldFormat{urldate}{\addcomma\space[\bibstring{urlseen}\space#1]}
\addto\captionsenglish{
  \renewcommand{\contentsname}
    {Table of Contents}
}

\begin{document}
\maketitle
\thispagestyle{empty}
\clearpage

\newpage
\pagenumbering{roman}
\thispagestyle{empty}
\tableofcontents
\clearpage

\newpage
\pagenumbering{arabic}

\section{Architecture requirements}
\subsection{Architectural scope}
\subsection{Quality requirements}
	\subsubsection{Performance}
		\begin{enumerate}
			\item \textbf{Description} \\
			\item \textbf{Justification} \\
			\item \textbf{Requirements}\\
		\end{enumerate}

	\subsubsection{Reliability}
		\begin{enumerate}
			\item \textbf{Description} \\
			\item \textbf{Justification} \\
			\item \textbf{Requirements}\\
		\end{enumerate}

	\subsubsection{Scalability}
		\begin{enumerate}
			\item \textbf{Description} \\
			\item \textbf{Justification} \\
			\item \textbf{Requirements}\\
		\end{enumerate}

	\subsubsection{Security}
		\begin{enumerate}
			\item \textbf{Description} \\
			\item \textbf{Justification} \\
			\item \textbf{Requirements}\\
		\end{enumerate}

	\subsubsection{Flexibility}
		\begin{enumerate}
			\item \textbf{Description} \\
			\item \textbf{Justification} \\
			\item \textbf{Requirements}\\
		\end{enumerate}

	\subsubsection{Maintainability}
		\begin{enumerate}
			\item \textbf{Description} \\
			\item \textbf{Justification} \\
			\item \textbf{Requirements}\\
		\end{enumerate}

	\subsubsection{Auditability/Monitorability}
		\begin{enumerate}
			\item \textbf{Description} \\
			\item \textbf{Justification} \\
			\item \textbf{Requirements}\\
		\end{enumerate}

	\subsubsection{Integrability}
		\begin{enumerate}
			\item \textbf{Description} \\
			\item \textbf{Justification} \\
			\item \textbf{Requirements}\\
		\end{enumerate}

	\subsubsection{Cost}
		\begin{enumerate}
			\item \textbf{Description} \\
			\item \textbf{Justification} \\
			\item \textbf{Requirements}\\
		\end{enumerate}

	\subsubsection{Usability}
		\begin{enumerate}
			\item \textbf{Description} \\
			\item \textbf{Justification} \\
			\item \textbf{Requirements}\\
		\end{enumerate}

\subsection{Integration and access channel requirements}
\subsection{Architectural constraints}

\section{Architectural patterns or styles}

\section{Architectural tactics or strategies}
 
\section{Reference architectures and frameworks}
 
\section{Access and integration channels}

\section{Technologies}
\subsection{Backend System}
\subsubsection{Programming Languages}
\subsubsection{Frameworks}
\subsubsection{Libraries}
\subsubsection{Database System}
\subsubsection{Operating System}
\subsubsection{Dependency Management and Build Tools}

\subsection{Web Interface}
\subsubsection{Programming Languages}
\subsubsection{Frameworks}
	\begin{itemize}
	\item AngularJS\\
		AngularJS is a JavaScript web application framework maintained by Google and volunteers, provding developers with a Model-View-Controller(MVC) and Model-View-ViewModel (MVVM) architectures to developer Single Page Applications (SPA). As the community has grown around AngularJS various plugins and components have been provided which eases the developement of rich internet applications, espesically with excellant support from AngularJS itself to support RESTful applications.

		The reason that AngularJS was chosen, is that it is a mature code base, with a large support community and which modern software engineering tools support to ensure a quality application is delivered to the end client. This allows one to do unit testing with ease, allow for the integration of a frontend end dependency management system such as bower and using modern build tools such as Grunt, in the end deliver a fast, stable, secure and efficient modern single page application.
	\end{itemize}

\subsubsection{Libraries}
	\begin{itemize}
		\item Modernizr\\
			Modernizr is a JavaScript(JS) library to provide developers a stable, consistent and standardized manner of feature detecion in browers. Based on the fact we will be targeting our application to modern HTML5 based browser, we need to utilize a feature detection toolkit in order to be able to gracefully degrade the website for users of older browsers.
		\item Karma\\
			Karma is a test runner for JavaScript development and being framework agnostic allows developers to select there own framework to be used with the testing. As we need to ensure that the system remains stable, we will be using Karma with Jasmine as the test framework for the frontend JavaScript unit testing.
		\item Jasmine\\
			Jasmine is a behavious driven testing framework for JavaScript with the advantage of not requiring further 3rd party libraries or needing a fully functional DOM to do the testing against.
		\item Gatling\\
			Gatling is a load testing framework used to test, measure, analyze and report on the performance of web services, allow one to ensure that a given project would be able to scale under increased load pressure.

			As scability is a requirement for this project, especially during month end peak load times, it is of the utmost importance to ensure that the system is scalable under this increased load. By using a testing framework like Gatling one will be able to measure quantitatively how scalable the system is.
	\end{itemize}

\subsubsection{Database System}
	\begin{itemize}
		\item Since we are utilizing a token based authenication mechanism, we need to store the access and refresh tokens on the client side, as the access token should accompany each request to the backend system.  The refresh token is also needed to refresh the access token, should the client continue to work longer than the period for which the access token is valid. We need to consider both newer and older client browsers, and hence we will utlize the following two persistance technologies:
		\begin{itemize}
			\item HTML5 Local Storage\\
				For clients who are utilizing an HTML 5 capable browser with local storage we will be using the newer and safer local storage ability provided by the be browser to store the token temporary. Upon logout of the backend system, these token will be earsed on the client side.
			\item Cookies\\
				To allow clients who are using and older browser to still use the system, we will store the token in cookies on the client system, which will be erased upon logout from the backend system, in order to increase security,
		\end{itemize}
	\end{itemize}

\subsubsection{Operating System}
	\begin{itemize}
		\item Device Operating System
			As we will be developing a SPA which will be running in the context of a web browser, the underlying device operating system is independant of the application, and hence we will be able to support any operating system which provides the required browser.
		\item Web Browser
			As we will be using AngularJS to build a SPA, we will be supporting the following browsers
		\begin{itemize}
			\item Google Chrome
			\item Mozilla Firefox
			\item Microsoft Internet Explorer 8,9 \& 10
			\item Microsoft Edge
			\item Apple Safari
		\end{itemize}
	\end{itemize}

\subsubsection{Dependency Management and Build Tools}
	\begin{itemize}
		\item Bower\\
			Bower is a package manager for frontend developement, allowing developers to manage HTML, JS, CSS and font 3rd party libraries.  It further allows for the installation of the correct versions of libraries with their corresponding dependencies. This allows developers to more cleanly determine on what libraries there projects depend and allows for better repository management. The reason is that most modern frontend libraries are already found in source version control repositories, and to add library source code to your own repository, not only inflates the repository size but doesn't allow for clean trakcing of the how the ibrary code itself has evolved over time.

		\item Grunt\\
			Grunt is a a JavaScript build tool that assist developers in automating reptitve tasks. Grunt being written in JavaScript makes the build system extremely flexable and allows the system to be adapted for any task required by developers, which has lead to the large amount of plugins being written for the build system.

			The reason for using grunt for the web development build is we require the following repetive tasks to be completed each time the web frontend needs deployment:
			\begin{itemize}
				\item Minification of CSS and JS source files.
				\item Concatenation and versioning in the filename of CSS and JS files.
				\item Removing of unneed CSS styles to decrease page loading time.
				\item Autowiring of Bower dependency libraries into the web frontend.
				\item Running of CSSLint and JSLint utilities on corresponding CSS and JS file to ensure the files are correct.
				\item Running of Karma JS unit tests to ensure JS functions are correct.
				\item Minification and compression on image files to ensure we decrease page loading time and increase user expierence.
				\item Running of Gatling stress testing to ensure changes made to frontend doesn't impact on the quality requirement for scalability and system performance.
			\end{itemize}
	\end{itemize}

\subsection{Android Client}
\subsubsection{Programming Languages}
	\begin{itemize}
		\item Java\\
			We will be writing a native Android application and hence we are required by the constraints of the Android system to utilize the Java programming language. In terms of Android using Java, Android relies heavy on the following Java concepts:
			\begin{itemize}
				\item Platform Independence\\
					After the developers have developed there applications in a selected programming language, the programmer often needs to compile there source code. The process of compiling source code is the conversion of the code written by the developer into a language that a certain device can understand. 
					
					It is important to note that different devices, understand different machine level langauges, and as such the developer is often forced to recompile the application for different devices understanding different machine languages. Using Java, Java source code is converted into a language called "bytecode", which is then interpreted and exectued by the Java Virtual Machine (JVM), which operates much like a physical CPU does.
					
					Using this ability of Java, it allows developers to compile there application once, and have it available on multiple different devices, by not needing to worry what machine language powers the given device. This gives developers the ability for the application to run on Android phones, tablets to even the latest toaster, car or water-bottle powered by Android.
					
				\item Secure Virtual Machines\\
					Since Java execute in a virtual machine, Android has taken advantage of this ability, and runs each application in its own virtual machine as seperate Linux users. In this way Android abstracts the execution of machine language away from the physical hardware, into a virtual machine which encapsulates and contains the bytecode and which can be more easily managed should applications not fulfill certain constraints, such as consuming too much power.
				\end{itemize}
		\item eXtensible Markup Language (XML)\\
			Android makes extensive use of XML to assist it in creating user screens in a process called inflation, storing string literals which allows for internationalization based on user preferences, and the selection of media such as image, video and audio based on the devices capablities, such as screen resolution, avaible decoders etc.
	\end{itemize}
	
\subsubsection{Frameworks}
	\begin{itemize}
		\item Spring for Android\\
		Spring for Android is a framework bringing selected capablities from the Spring framework to the Android platform, allowing for faster, cleaner and more efficent developement of Andriod applications. Capabilities of Spring for Android include:
		\begin{itemize}
			\item A RESTful client with full Hypermedia as the Engine of Application State (HATEOAS) capability
			\item Authentication support which includes OAuth 1.0 and 2.0 capability, allow for token based authentication.
		\end{itemize}
	\end{itemize}

\subsubsection{Libraries}
	\begin{itemize}
		\item Android Butterknife\\
		Butterknife for Android brings Dependency Injection (DI) capability to Android's view system, allowing for cleaner and faster code. This is achieved by using Java annotations to bind visual components to corresponding Android Java components, thereby eliminating the need for poluting production source code with the infamous Android findViewById() method, which itself uses slow Java reflection to accomplish its task. Butterknife on the other hand produces code at compile time to assist in the lookup process.
	\end{itemize}

\subsubsection{Database System}
	\begin{itemize}
		\item Couchbase Mobile\\
		Couchbase Mobile is a NoSQL database system supporting all major mobile platforms, providing fast and consistent access to data, with or without a network connection. The Couchbase Mobile system consists of three key components:
		\begin{itemize}
			\item Couchbase Lite\\
			A mobile orientated NoSQL database allowing for low latency read/write access of data.
			\item Couchbase Sync Gateway\\
			An internet facing component, interfacing between mobile clients and the backend system to achieve secure synchronization of data.
			\item Couchbase Server\\
			A scalable, high performance, cluster based and enterprise based NoSQL database server providing NoSQL database functionality, with traditional based Structed Query Language functionality.
		\end{itemize}

		\item Utilizing a complete turn-around database system like Couchbase Mobile, will allow for the delivery of a fast and consistent user expierence, with the additional benefit of providing the mobile client with offline access and capability to view and edit selected data. 
	\end{itemize}

\subsubsection{Operating System}
\subsubsection{Dependency Management and Build Tools}



\end{document}
